\section{Planitud}
\label{section:0.6-fr}

\begin{env}[6.0]
\label{0.6.0.0-fr}
La noción de planitud se debe a J.-P.~Serre [16]; en lo que sigue, se omiten
las demostraciones de los resultados que se exponen en el \emph{Álgebra
conmutativa} de N.~Bourbaki, a la que remitimos al lector. Suponemos que
todos los anillos son conmutativos\footnote{Véase la exposición citada de
N.~Bourbaki para la generalización de la mayoría de los resultados al caso no
conmutativo.}.

\oldpage[0\textsubscript{I}]{55}

Si $M$, $N$ son dos $A$-módulos, $M'$ (resp. $N'$) un submódulo de $M$
(resp. $N$), denotaremos por $\operatorname{Im}(M'\otimes_A N')$ el submódulo
de $M\otimes_A N$, imagen de la aplicación canónica
$M'\otimes_A N'\to M\otimes_A N$.
\end{env}

